\section{KPI và cây phân tích}\label{sec:kpi}
Bảng \ref{tab:kpi} liệt kê các KPI chính; mỗi KPI được định nghĩa một lần bằng công thức DAX (Phụ lục A) và đối chiếu với giá trị tính bằng SQL trên kho. \textbf{GMV} chỉ tính đơn không thuộc nhóm Failed (hủy hoặc không khả dụng), vì vậy nhỏ hơn tổng giá hàng của mọi đơn (13,49 triệu so với 13,59 triệu BRL). KPI \emph{ngày hứa giao so với ngày giao thực tế} so sánh số ngày Olist hứa với số ngày giao thực tế: ngày dự kiến rất rộng rãi (trung bình 24,4 ngày) trong khi giao thực tế trung bình 12,6 ngày, nên tỷ lệ ``đúng hạn'' cao (93,2\%) chưa phản ánh trải nghiệm.
\begin{table}[htbp]
\centering\small
\caption{Các chỉ số KPI chính (tính bằng SQL trên kho; đối chiếu với Power BI)}
\label{tab:kpi}
\begin{tabular}{lrr}
\toprule
\textbf{KPI} & \textbf{Toàn bộ dữ liệu} & \textbf{Cửa sổ 2017-01 → 2018-08} \\
\midrule
GMV & 13.494.400,74 & 13.449.529,68 \\
Orders & 98.199 & 97.905 \\
AOV & 137,42 & 137,37 \\
Freight ratio & 16,6\% & 16,6\% \\
Orders placed & 99.441 & 99.092 \\
Cancel rate & 1,24\% & 1,19\% \\
On-time rate & 93,2\% & 93,2\% \\
Avg lead days & 12,56 & 12,54 \\
Avg promised days & 24,37 & 24,28 \\
Promised vs actual (days) & 11,81 & 11,75 \\
Avg review & 4,09 & 4,09 \\
\% negative & 14,7\% & 14,6\% \\
Repeat customer rate & 3,12\% & 3,11\% \\
MQLs (theo tháng liên hệ đầu) & 8.000 & 8.000 \\
Won deals & 842 & 842 \\
MQL to won rate & 10,5\% & 10,5\% \\
Avg days to close & 48,50 & 48,50 \\
Sellers with first sale & 380 & 380 \\
Won to first sale rate & 45,1\% & 45,1\% \\
\bottomrule
\end{tabular}
\\[2pt]{\footnotesize GMV chỉ tính đơn không thuộc nhóm Failed (hủy / không khả dụng); tổng giá hàng mọi đơn là 13.591.643,70 BRL.}
\end{table}

Mỗi dashboard có một cây phân tích đi từ mục tiêu đến câu hỏi, chỉ số và chiều (Hình \ref{fig:cay-d1} đến \ref{fig:cay-d7}).
\hinh{3-2-7_cay_phan_tich_D1.png}{Cây phân tích dashboard D1 -- Executive Overview}{fig:cay-d1}{0.75}
\hinh{3-2-7_cay_phan_tich_D2.png}{Cây phân tích dashboard D2 -- Sales và Product}{fig:cay-d2}{0.75}
\hinh{3-2-7_cay_phan_tich_D3.png}{Cây phân tích dashboard D3 -- Order Fulfillment}{fig:cay-d3}{0.75}
\hinh{3-2-7_cay_phan_tich_D4.png}{Cây phân tích dashboard D4 -- Geography và Freight}{fig:cay-d4}{0.75}
\hinh{3-2-7_cay_phan_tich_D5.png}{Cây phân tích dashboard D5 -- Customer Experience}{fig:cay-d5}{0.75}
\hinh{3-2-7_cay_phan_tich_D6.png}{Cây phân tích dashboard D6 -- Seller Acquisition}{fig:cay-d6}{0.75}
\hinh{3-2-7_cay_phan_tich_D7.png}{Cây phân tích dashboard D7 -- Data Quality và ETL Monitor}{fig:cay-d7}{0.75}
\todo{Nhóm vẽ lại các cây phân tích bằng draw.io nếu muốn thống nhất kiểu trình bày; thêm bảng KPI đầy đủ theo định nghĩa trong docs/DW\_DESIGN.md mục 9 (công thức, nguồn, đơn vị).}
